\title{Byte Latent Transformer: Patches Scale Better Than Tokens}

\begin{abstract}
We propose the Byte Latent Transformer ({\textsc BLT}), a novel large language model architecture that operates at the byte level, achieving for the first time parity with token-based LLMs in terms of performance at scale, while delivering marked gains in inference speed and reliability. {\textsc BLT} processes inputs as dynamically determined byte patches, treating these as fundamental computational units. The formation of patches is governed by the entropy of the upcoming byte, allowing the model to assign additional resources and representational power to segments with greater informational complexity. We conduct and report the first scaling analysis under controlled floating-point operation constraints for byte-level LLMs, reaching up to 8 billion parameters and utilizing 4 trillion training bytes. Our findings establish that it is practical to train models directly on raw byte sequences without predefined vocabularies, with improvements observed in both training and inference efficiency due to adaptively opting for longer patches in more predictable contexts. Additional enhancements are evident in the model’s reasoning abilities and its handling of rare or complex scenarios. Taken together, {\textsc BLT} attains superior scaling properties compared to tokenization-based architectures for equivalent inference budgets, facilitated by joint scaling of both patch length and model size.
\end{abstract}

\section{Introduction}
\label{section:intro}

We present the Byte Latent Transformer~(\textbf{BLT}), an architecture that eliminates the need for tokenization by directly learning from raw byte sequences. This approach, for the first time, achieves parity with models that rely on tokenization when scaled up, while also delivering notable gains in computational efficiency and resilience (\S\ref{sec:robustness}). In contrast, current large language models (llms) are predominantly trained in an end-to-end manner except for the crucial preliminary step of tokenization, which heuristically segments byte streams into a fixed vocabulary of tokens. This process introduces inherent biases in string compression, resulting in various limitations including sensitivity to particular domains or modalities~\citep{dagangetting}, increased vulnerability to noisy inputs~(\S\ref{sec:robustness}), insufficient handling of orthography~\citep{edman2024cute}, and disparities across languages~\citep{liang2023xlm,petrov2024language,limisiewicz2024myte}.

Tokenization has traditionally played a critical role, as training large language models on raw byte sequences incurs substantial computational expense due to extended sequence lengths~\citep{xue2022byt5}. Some studies address this challenge by introducing more efficient variants of self-attention~\citep{el2020characterbert, clark2022canine} or by utilizing architectures that dispense with attention mechanisms altogether~\citep{wang2024mambabyte}~(\S\ref{section:related_work}). Nevertheless, such improvements are mainly beneficial for training \textit{smaller models}. When scaling up, the majority of computational resources are consumed by the sizeable feed-forward network layers operating on each byte, rather than by the attention modules themselves.

For effective compute allocation, we introduce a dynamic, trainable approach to clustering bytes into \textit{patches}~(\S\ref{section:patching}), alongside a novel model architecture that combines both byte-level and patch-level data. In contrast to traditional tokenization, {\textsc BLT} employs no predetermined patch vocabulary. Instead, any selection of bytes can be converted into latent patch embeddings using lightweight, learned encoder and decoder mechanisms. Our findings demonstrate that this strategy allows for \textit{greater} compute efficiency compared to models that rely on tokenization.

Tokenization-based language models distribute an equal amount of computational effort to every token. This uniform treatment sacrifices efficiency for the sake of performance, as the underlying compression heuristics driving token induction do not always align with the actual difficulty of the predictions. A key innovation in our design is the principle that models ought to vary compute allocation in response to the complexity of the task at hand. For instance, most word endings are relatively straightforward to generate and involve decisions with lower entropy compared to, say, selecting the first word of a sentence, making extensive transformer resources unnecessary in these cases. {\textsc BLT}'s framework (\S\ref{section:architecture}) embodies this idea by incorporating three transformer blocks: a pair of lightweight byte-level \textit{local models} complemented by a more substantial global \textit{latent transformer}~(\autoref{fig:arch_high_level}).

In order to decide how bytes should be aggregated into patches—and thus, how to judiciously assign compute—{\textsc BLT} partitions inputs according to the entropy associated with the prediction of the following byte, resulting in groups of bytes that share comparable levels of information density.

We conduct the inaugural scaling analysis of byte-level models managed by flop budgets, extending up to 8B parameters and trained on 4T bytes, thereby demonstrating the feasibility of training large-scale models directly from raw byte sequences without the constraints of pre-defined tokenization schemes. Crucially, {\textsc BLT} attains comparable performance to Llama 3 under matched training computational budgets\footnote{The computation cost for a model is quantified by tallying the total number of Floating Point OPerations (flops) expended.}, while reducing inference flops by as much as 50\%~(\S\ref{section:scaling}). 
Furthermore, our study finds that direct byte-level modeling substantially enhances performance on the infrequent and idiosyncratic patterns in the data. {\textsc BLT} exhibits increased robustness to input noise relative to models relying on tokenizers, and displays improved abilities in character-level reasoning—achievements evidenced through experiments in orthography, phonological tasks, and low-resource machine translation~(\S\ref{sec:robustness}).
Additionally, with the {\textsc BLT} framework, we demonstrate the capacity to scale both the model and the patch sizes simultaneously without exceeding the same inference flops ceiling. Employing larger patch sizes generally leads to computational savings; these savings can then be directed towards enlarging the global latent transformer, which subsequently requires fewer updates. Our inference-flop constrained scaling experiments~(\autoref{fig:fixed_inference_scaling}) highlight marked improvements in scaling dynamics when compared to traditional tokenization-driven approaches.

To conclude, this work offers several key contributions: 1) We present {\textsc BLT}, a byte-level latent LLM framework that adaptively distributes computational resources to boost flop efficiency; 2) We demonstrate that our approach matches Llama 3’s training flop-controlled performance at scales up to 8B parameters, while also permitting significant increases in flop efficiency—up to 50\%—with only minimal reductions in standard evaluation scores; 3) {\textsc BLT} enables a novel scaling strategy for LLMs, allowing model capacity to grow without exceeding a constant inference compute ceiling; 4) Our results highlight {\textsc BLT}’s enhanced resilience to input perturbations and its sensitivity to sub-word data features typically overlooked by token-based LLMs. The code for training and inference with {\textsc BLT} is made publicly available at~\url{https://github.com/facebookresearch/blt}.

\section{Patching: From Individual Bytes to Groups of Bytes}
\label{section:patching}

Dividing bytes into \textit{patches} enables {\textsc BLT} to flexibly adjust computational resources in response to contextual demands. As illustrated in Figure~\ref{fig:patching_types}, multiple strategies exist for grouping bytes into patches. More formally, a patching function $f_p$ partitions a byte sequence $\pmb{x} = \{x_i, |i=1,\ldots, n\}$ of length $n$ into $m < n$ patches $\pmb{p} = \{p_j | j = 1, \ldots, m\}$ by associating each $x_i$ with the binary set \{0,1\}, where a 1 denotes the beginning of a new patch. In both token-centric and patch-oriented architectures, the principal determinant of computational expenditure is the number of operations performed by the core Transformer; within {\textsc BLT}, this corresponds to the number of patches generated for a particular patching method. Thus, the average number of bytes grouped within each patch—referred to as the \textit{patch size}—becomes the primary metric for assessing processing cost during both model training and evaluation using a specific patching technique~(\S\ref{section:flops}).

In the subsequent sections, we detail three distinct patching mechanisms: segmenting by a constant byte count per patch~(\S\ref{section:static-patch}), patching at whitespace boundaries~(\S\ref{section:white-patch}), and dynamically forming patches by leveraging byte-level entropy estimates from a compact language model~(\S\ref{section:dyn-patch}). We conclude by exploring incremental patching approaches and delineating the differences between patching and tokenization~(\S\ref{section:bpe-patch}).

\subsection{Strided Patching Every K Bytes}
\label{section:static-patch}

A simple approach to organizing bytes involves segmenting them into fixed-size patches of length $k$, as in MegaByte~\citep{yu2023megabyte}. Using a constant stride simplifies both the implementation of training and inference processes, while also providing an intuitive way to adjust average patch length and directly manage computational complexity. Despite these practical benefits, this patching strategy presents notable limitations. Primarily, computational resources are not distributed according to data complexity: for instance, a transformer step $j$ might be inefficiently applied to predict whitespace in source code or, conversely, underutilized for information-rich segments like mathematical notation. Additionally, this approach can yield inconsistent and context-insensitive segmentations of analogous byte sequences; for example, the same word may be partitioned in divergent ways.

\subsection{Space Patching}
\label{section:white-patch}

\citet{slagle2024spacebyte} introduces a straightforward yet powerful modification to strided patching by generating new patches at every space-like byte\footnote{
    Space-like bytes are considered any byte except latin characters, digits, or utf-8 continuation bytes. Moreover, it is required that every patch must include at least one byte that is not space-like.}, which typically serve as natural demarcations between linguistic elements across numerous languages. In Space patching, a separate latent transformer operation—meaning increased computational effort—is applied to each word, enabling words to be partitioned in a consistent manner regardless of their position in a sequence. This approach also strategically allocates computational resources to challenging prediction points that frequently occur after spaces. For instance, in answering ``Who composed the Magic Flute?\uline{\hspace{.6em}}'', predicting the very first byte of the correct answer is significantly more complex than predicting the subsequent bytes, as that initial character narrows the possible candidates dramatically, thereby rendering the remainder of the answer, such as ``Mozart'', much simpler to complete. Nevertheless, space patching exhibits limitations: it cannot efficiently accommodate all languages and domains, and is restricted by its inability to adapt the patch size. In what follows, we present an alternative patching strategy, which builds upon the observation that the initial bytes of words are usually the most challenging to forecast, while also permitting dynamic control over patch dimensions.

\subsection{Entropy Patching: Using Next-Byte Entropies from a Small Byte LM}
\label{section:dyn-patch}

In lieu of using rule-based heuristics like whitespace, our method adopts a data-centric strategy for pinpointing regions with high uncertainty in next-byte predictions. We present \textit{entropy patching}, a technique that establishes patch boundaries through the analysis of entropy values.

To this end, we train a compact byte-level auto-regressive language model using the {\textsc BLT} training dataset. Next, we calculate the next byte entropies from the language model's probability distribution $p_{e}$ over the byte vocabulary $\mathcal{V}$:

\begin{equation}
H(x_i) = \sum_{v \in \mathcal{V}} p_{e}(x_i=v|\pmb{x}_{<i}) \log p_{e}(x_i=v|\pmb{x}_{<i})
\end{equation}

We explore two approaches for detecting patch boundaries based on the entropies $H(x_i)$. The initial method involves selecting points where the entropy surpasses a set global threshold, as depicted in~\autoref{fig:patching}. The alternative approach highlights points whose entropies are significantly higher compared to the immediately preceding value. This latter method may also be viewed as pinpointing locations where the expected roughly monotonic decline in entropy within a patch is disrupted.

\begin{align*}
    \text{Global Constraint}\hspace{1cm}H(x_t)& > \theta_g\\
    \text{Approx. Monotonic Constraint}\hspace{1cm}H(x_t)& - H(x_{t-1}) > \theta_r
\end{align*}

Patch boundaries are determined in a computationally inexpensive preprocessing phase that takes place while data is being loaded. This approach contrasts with~\citet{nawrot-etal-2023-efficient}, who employ a classifier to estimate patch boundaries based on entropy predictions. In our experimental analysis~(\S\ref{section:experiments}), we evaluate and contrast these two techniques for separating low entropy bytes from high entropy ones.

\subsection{The Byte-Pair Encoding (BPE) Tokenizer and Incremental Patching}
\label{section:incremental-patching}
\label{section:bpe-patch}

A number of state-of-the-art llms, such as our reference model Llama 3, rely on a subword tokenization strategy like bpe~\citep{gage1994new, sennrich-etal-2016-neural}. In this context, we define ``tokens'' as groups of bytes selected from a \textit{finite} vocabulary that is set before model training, whereas ``patches'' refer to dynamically formed sequences that are not constrained by a fixed vocabulary. An essential distinction is that models using tokens lack direct access to the raw byte-level information, which sets them apart from approaches utilizing patches.

An important advancement offered by {\textsc BLT} relative to traditional tokenization-based systems is its novel adjustment of the balance between vocabulary size and computational resources. In conventional large language models, expanding the vocabulary results in tokens that cover more text on average, thereby reducing the number of decoding steps. However, this also necessitates a larger output dimension in the model's final projection layer. Consequently, tokenization-based methods are constrained in how much they can vary both the token length and the associated computational cost during inference. As an illustrative case, Llama 3 raises its average token length from 3.7 to 4.4 bytes, but this comes at the expense of enlarging its embedding table to four times the size used in Llama 2.

During generation, {\textsc BLT} must determine if the present position in the byte sequence coincides with a patch boundary, as this influences whether the Latent Transformer will be utilized for additional computation. This assessment must be made without knowledge of bytes that are yet to be produced; in other words, the process for partitioning the sequence into patches cannot rely on access to subsequent bytes. To formalize, a patching method $f_p$ is said to possess incremental patching if the following holds:
$$f_p(\pmb{x}_{<i}) = f_p(\pmb{x})_{<i}$$
The bpe approach does not exhibit incremental patching properties, since a given prefix may yield varying tokenizations based on what tokens come next, and thus it fails to fulfill the above requirement\footnote{Incorporating a dedicated delimiter token to mark patch boundaries can enable bpe to function incrementally, but at the cost of lengthening the byte sequence.}.

\section{{\textsc BLT} Architecture}
\label{section:architecture}
{\textsc BLT} consists of a large-scale global autoregressive language model that processes patch representations, supplemented by two more compact local models: one to encode sequences of bytes into patches, and another to reconstruct bytes from patch representations~(\autoref{fig:arch_high_level}).

\subsection{Latent Global Transformer Model}

\textit{The Latent Global Transformer} refers to an autoregressive transformer, denoted as $\mathcal{G}$, comprising $l_{\mathcal{G}}$ layers, which takes a sequence of latent input patch representations, $p_j$, and produces a corresponding sequence of output patch representations, $o_j$. Throughout this work, the subscript $j$ identifies individual patches, while $i$ indicates bytes. The architecture incorporates a block-causal attention mask~\citep{dubey2024llama}, ensuring that attention is limited to patches up to and including the current one within the same document. Since the global model is responsible for the majority of floating-point operations both during pre-training and at inference time, controlling when the model is activated enables the modulation and allocation of computational resources according to the varying complexity of different input and output segments.

\subsection{Local Encoder}
\label{section:local}

\textit{The Local Encoder Model}, represented as $\mathcal{E}$, is designed as a compact transformer-based architecture consisting of $l_{\mathcal{E}} << l_{\mathcal{G}}$ layers. Its primary objective is to rapidly convert an input sequence of bytes $b_i$ into informative patch embeddings, $p_j$. Notably, unlike standard transformers, this model introduces a cross-attention mechanism after each transformer layer, which serves to aggregate byte-level information into patch-level representations~(\autoref{fig:crossattn}).

Initially, the input bytes $b_i$ are mapped into a high-dimensional space using an embedding matrix of size $\mathbb{R}^{256 \times h_{\mathcal{E}}}$, resulting in embedded vectors $x_i$. Optionally, these embeddings can be enriched with hash-embeddings to provide supplementary information~(\S\ref{sec:hash-emb}).

Through a sequence of alternating transformer and cross-attention layers~(\S\ref{sec:enc-cross-attn}), these embedded byte representations are transformed into patch representations $p_i$. These patches serve as input to the global transformer component, $\mathcal{G}$. The attention mechanism within the transformer employs a \textit{local block causal} mask, such that each byte can only attend to a predefined window of $w_{\mathcal{E}}$ previous bytes. This attention may extend across dynamically set patch boundaries but is restricted from spanning document boundaries. The upcoming subsections elaborate on the specifics of the embedding strategy and the configuration of the cross-attention module.

\subsubsection{Encoder Hash n-gram Embeddings}
\label{sec:ngram-emb}
\label{sec:hash-emb}
An essential aspect of building strong and informative representations for each position $i$ involves encoding the contextual information from the bytes that come before. In {\textsc BLT}, this is accomplished by treating the byte $b_i$ not only in isolation but also within the context of its surrounding byte n-gram. At every step $i$, the process begins with the creation of byte-grams

\begin{align}
    g_{i,n}=\{b_{i-n + 1},\ldots, b_i\}
\end{align}

for each byte position $i$ and $n$ from three to eight.\footnote{
We omit byte-grams of size $n$ or more when $i<n$. }

Next, we present hash $n$-gram embeddings, which assign each byte $n$-gram to a position in a fixed-size embedding table $E_{n}^{hash}$ by utilizing a hash function, for every $n$ value in $\{3,4,5,6,7,8\}$~\citep{bai2010learning}. The computed $n$-gram embedding is subsequently summed with the embedding of the individual byte. This combined representation is then normalized and supplied as input to the local encoder model. The calculation of the augmented embedding

\begin{align}
e_i &= x_i + \sum_{n = 3,...,8}E_{n}^{hash}(\text{Hash}(g_{i,n})) \\
\text{where, Hash}(g_{i,n}) &= \text{RollPolyHash}(g_{i,n}) \% |E_{n}^{hash}|
\end{align}

We scale $e_i$ by dividing it by the sum of the total $n$-gram sizes and one, applying the RollPolyHash method as detailed in Appendix~\ref{appendix:rollpolyhash}. Section~\ref{section:ablations} presents an analysis of how varying $n$-gram hash embedding parameters—including the size of $n$ and the dimensions of the embedding table—impacts trends observed in flop-controlled scaling laws. Alongside the hash-based $n$-gram embeddings, we explored frequency-based $n$-gram embeddings, with a comprehensive discussion provided in Appendix~\ref{appendix:ngrams}.

\subsubsection{Encoder Multi-Headed Cross-Attention}
\label{sec:enc-cross-attn}
Our approach generally mirrors the structure of the input cross-attention module from the Perceiver architecture~\citep{jaegle2021perceiver}, with the principal modification that latent variables now represent the embeddings of individual, dynamic patches instead of a fixed group of latent vectors~(\autoref{fig:crossattn}). Each patch, in this setup, is only cross-attended with the bytes it contains. Specifically, for each patch $p_j$, we construct a query by pooling the byte-level representations within $p_j$ and subsequently applying a linear transformation via $\mathcal{E}_{C} \in \mathbb{R}^{h_{\mathcal{E}} \times (h_{\mathcal{E}} \times U_{\mathcal{E}})}$, where $U_{\mathcal{E}}$ denotes the number of cross-attention heads used in the encoder. Explicitly, for the bytes belonging to patch $p_j$ (denoted as $f_{\text{bytes}}(p_j)$), this operation is formalized as follows:

\begin{align}
P_{0,j} &= \mathcal{E}_{C}(f_\text{bytes}((p_j)), f~ \text{is a pooling function} \\
P_l &= P_{l-1} + W_o\left(\text{softmax}\left(\frac{QK^T}{\sqrt{d_k}}\right)V\right) \\
\text{where } Q_j &= W_q(P_{l-1,j}), K_i = W_k(h_{l-1,i}), V_i = W_v(h_{l-1,i}) \\
h_l &= \text{Encoder-Transformer-Layer}_l(h_{l-1})
\end{align}

Here, $P \in \mathbb{R}^{n_p \times h_{\mathcal{G}}}$ denotes $n_p$ patch-level vectors, which are passed through the global module. These vectors are initially obtained by pooling the byte-level embeddings $e_i$ assigned to each patch $p_j$. The matrices $W_q$, $W_k$, $W_v$, and $W_o$ serve as the linear transformation weights for computing queries, keys, values, and outputs, respectively. The keys and values are derived by projecting the byte encodings $h_i$ generated from the preceding layer (for the initial layer, these are $e_i$). We adopt a patch-specific masking approach such that for any given query $Q_j$, only the keys and values linked to the bytes within patch $j$ are attended to. Owing to the use of multi-head attention across $Q, K$, and $V$, and given that patch vectors generally possess greater dimensionality ($h_{\mathcal{G}}$) than $h_{\mathcal{E}}$, we represent $P_l$ using multiple heads, each of dimension $h_{\mathcal{E}}$ during cross-attention; these are subsequently concatenated to match the $h_{\mathcal{G}}$ dimensionality. We further apply a pre-LayerNorm operation to queries, keys, and values. Positional embeddings are intentionally omitted from this cross-attention design. Lastly, the cross-attention module includes a residual connection encasing the entire block.

\subsection{Local Decoder} 

In parallel with the local encoder, the local decoder $\mathcal{D}$ is also designed as a lightweight transformer, comprising significantly fewer layers than the global model, specifically $l_{\mathcal{D}} << l_{\mathcal{G}}$. This component takes a sequence of global patch embeddings, denoted as $o_j$, and reconstructs them into a sequence of raw bytes, $y_i$. The generation of raw bytes is conditioned on the sequence of bytes previously decoded, using as input the hidden states output by the local encoder for the byte sequence. The architecture consists of $l_{\mathcal{D}}$ stacked blocks, each alternating between cross-attention mechanisms and transformer layers. Within each block, cross-attention is applied first to bridge patch representations and byte representations, following which the transformer layer further processes the derived byte-level representations.

\subsubsection{Decoder Multi-headed Cross-Attention} 

Within the cross-attention mechanism of the decoder, the query and key/value assignments are inverted compared to the standard setting: here, the queries correspond to the byte representations, while the patch representations serve as the keys and values. The byte representations that serve as the starting point for cross-attention are taken directly from the byte embeddings produced by the final encoder layer, denoted as $h_{l_{\mathcal{E}}}$. At layer $l$, the byte representation $d_{l,i}$ is then updated as follows:

\begin{align}
D_0 &= h_{l_{\mathcal{E}}} \\
B_l &= D_{l-1} + W_o\left(\text{softmax}\left(\frac{QK^T}{\sqrt{d_k}}\right)V\right), \\
  &\text{where } Q_i = W_q(d_{l-1,i}),\;\; K_i = W_k(\mathcal{D}_C(o_j)),\;\; V_i = W_v(\mathcal{D}_C(o_j)) \\
  D_l &= \text{Decoder-Transformer-layer}_l(B_l)
\end{align}

Here, as before, $W_k$ and $W_v$ denote the key and value projection matrices, respectively, both of which conduct linear transformations and splitting via $\mathcal{D}_C$ on the final patch embeddings $o_j$ produced by the global model. The projection matrix $W_q$ is responsible for projecting queries from the byte-level representations $d_{l-1}$ that come from the preceding decoder transformer layer (or from $h_{l_{\mathcal{E}}}$ in the case of the initial layer). The output is subsequently projected with $W_o$. Consequently, $B$ resides in $\mathbb{R}^{h_{\mathcal{D}} \times n_b}$, where $n_b$ indicates the number of output bytes. The subsequent decoder hidden states, $D_l$, are then obtained by applying a decoder transformer layer to the result of the cross-attention computation, $B$. Mirroring the approach used in the local encoder's cross-attention, we deploy several attention heads, implement pre-LayerNorms, forego positional embeddings, and incorporate a residual pathway wrapping the cross-attention mechanism.

\section{Experimental Setup}
\label{section:experiments}
We devised rigorous and standardized experiments that systematically compare {\textsc BLT} and models utilizing tokenization, paying careful attention to control for sequence length in order to avoid inadvertently favoring {\textsc BLT} by allowing it access to longer contexts.

\subsection{Pre-training Datasets}  
The models evaluated in this study are pre-trained on two main corpora: 1) The Llama 2 dataset~\citep{touvron2023llama}, consisting of 2 trillion tokens sourced from diverse, publicly available datasets, which are then subject to extensive cleaning and filtering processes to enhance their overall quality; and 2) {\textsc BLT}-1T, a novel collection comprising 1 trillion tokens obtained from various public domains, further supplemented with a portion of the pre-training data provided by Datacomp-LM~\citep{li2024datacomp}. The Llama 2 dataset is utilized primarily for scaling law analyses on the optimal token quantity, as previously characterized by~\cite{dubey2024llama}, in order to inform key architectural decisions for {\textsc BLT}. In contrast, the {\textsc BLT}-1T dataset serves for full pre-training runs, enabling comparative evaluations against Llama 3 on downstream performance benchmarks. Importantly, both datasets are devoid of any content derived from Meta platforms or services. For the baseline comparisons involving models utilizing tokenizers, we adopt the Llama 3 tokenizer, which features a 128K token vocabulary; this choice was motivated by its consistently stronger baseline performance relative to the Llama 2 tokenizer in our empirical assessments.

\subsection{Entropy Model}  
For our experiments, the entropy model consists of a byte-level language model that is trained using the same data distribution as the complete {\textsc BLT} model. Unless specified otherwise, this component uses a transformer architecture with 100M parameters, consisting of 14 layers, each with a hidden dimension of 512, and employs sliding window attention spanning 512 bytes. All other hyperparameters align with those used in both our local and global transformer settings. Various architectures, model capacities, and receptive fields were explored as outlined in \autoref{section:ablations}. Notably, if the model's receptive field is sufficiently limited, the resulting trained entropy model can be represented effectively using a compact lookup table.

\subsection{Entropy Threshold and Equalizing Context Length} For models employing entropy-driven patching strategies, we determine an appropriate patching threshold to yield a specific average \textit{patch size} across the pretraining data mixture. In {\textsc BLT}, in contrast to conventional tokenization, the \textit{patch size} is a freely adjustable parameter, which considerably affects the effective context length available to the model. To preserve comparable average context lengths and avoid conferring an artificial advantage to models processing larger patches, we hold the total number of bytes per batch constant in expectation. Consequently, for configurations using larger patch sizes, the sequence length is reduced correspondingly. When training on Llama 2 data, we set the context to 8k bytes, whereas for the {\textsc BLT}-1T dataset, we increase the context window to an average of 16k bytes, all while keeping the mean batch size fixed at 16M bytes.

Although the mean batch size does not vary, dynamic patching techniques produce varying proportions of bytes relative to patches in each batch. In our {\textsc BLT} training framework, we optimize performance by assembling batches based on the number of patches, which eliminates the need for padding within the computationally intensive latent transformer. This approach guarantees a fixed patch count per batch. To prevent memory issues that could arise from sequences containing exceptionally large patches during training, we pad and, if needed, truncate byte sequences to lengths of 12k and 24k for the Llama 2 and {\textsc BLT}-1T datasets, respectively.

\subsection{Entropy Model Context}
\label{section:entropy-context}
Through our empirical analysis, we observe that entropy patching tends to result in increasingly larger patches within structured formats, such as multiple choice tasks, which often feature substantial repetition (see the example from MMLU in \autoref{fig:mmlu_patching}). This phenomenon can be attributed to the reduced entropy associated with recurring elements in the entropy model context. Consequently, for the full-scale deployment of {\textsc BLT}-Entropy at patch size 4.5, we refresh the entropy context by introducing new lines and implement an approximate monotonicity constraint, which mitigates the issue of "entropy drift" arising from fluctuations in context length. Notably, this modification solely influences our entropy calculations; the procedure for determining the appropriate entropy threshold remains unchanged.

\subsection{FLOPs Estimation} 
\label{section:flops}

Our methodology for calculating transformer flops is principally based on the equations presented in Chinchilla~\citep{hoffmann2022training}, accounting for the computational costs in the feed-forward components, the qkvo projections within self-attention, as well as attention calculation and output projection. Unlike previous approaches, we treat the input embedding stage as a highly efficient lookup rather than relying on dense matrix operations, effectively attributing it zero flops. In line with prior studies, we approximate the cost of the backward pass as being double that of the forward pass.

To determine the flops \textit{per byte} metric for {\textsc BLT} models, we aggregate the computational costs from several components: the local encoder transformer, the global latent transformer, the local decoder transformer, as well as the cross-attention mechanisms within both the encoder and decoder. This results in the following expression:

\begin{align}
    \text{FL}_{\text{{\textsc BLT}}} &= \frac{\text{Transf. FL}(h_{\mathcal{G}}, l_{\mathcal{G}}, m=n_{ctx}/n_p, V=0)}{n_p}\\
   &+\text{Transf. FL}(h_{\mathcal{E}}, l_{\mathcal{E}}, m=w_{\mathcal{E}}, V=0) \\
   &+ \text{Transf. FL}(h_{\mathcal{D}}, l_{\mathcal{D}}, m=w_{\mathcal{D}}, V=256) \\ 
    &+ \text{Cross Attn. FL}(h_{\mathcal{E}}, l_{\mathcal{E}}, m=n_p, r=n_p/k) \times \frac{k}{n_p}\\ 
   &+ \text{Cross Attn. FL}(h_{\mathcal{D}}, l_{\mathcal{D}}, m=k, r=k/n_p) 
\end{align}

In this context, $n_{ctx}$ denotes the length of the sequence measured in bytes, while $n_p$ indicates the patch size. The parameter $r$ specifies the proportion of queries relative to keys and values, and $k$ refers to the ratio between the patch and byte dimensions—that is, it represents the count of separate local model divisions merged to yield the comprehensive model representation (with $k=2$ as illustrated in \autoref{fig:crossattn}). The term $V$ stands for the vocabulary size used during the output projection, which appears exclusively in the local decoder. The choice of context length $m$ for the attention mechanism depends on whether the operation is performed over the byte or the patch sequence, leading to adjustments in the computation of attention. We adapt the formulas for attention flops for each module accordingly. Detailed expressions for the computation of Transformer-FLOPs and Cross-Attention FLOPs can be found in Appendix~\ref{section:flops-appendix}.

\subsection{Bits-Per-Byte Estimation} Since perplexity is defined relative to a specific tokenizer and quantifies uncertainty at the token level, it becomes less meaningful when models employ different tokenization schemes. In line with prior research~\citep{xue2022byt5, yu2023megabyte, wang2024mambabyte}, we utilize Bits-Per-Byte (BPB) as an alternative metric. BPB represents a tokenizer-agnostic analog to perplexity, enabling fairer comparison between models operating at the byte and token levels. Concretely:

\begin{align}
    \text{BPB}(x)&= \frac{\mathcal{L}_{CE}(\pmb{x})}{\ln(2)\cdot n_{\text{bytes}}}
\end{align}

where the uncertainty over the data $\pmb{x}$ as measured by the sum of the cross-entropy loss is normalized by the total number of bytes in $\pmb{x}$ and a constant. 

\subsection{Transformer Architecture Hyperparameters} In configuring the transformer blocks within {\textsc BLT}, encompassing both the local and global model variants, we adopt design choices primarily in line with the Llama~3 framework~\citep{dubey2024llama}. Specifically, our feed-forward sublayers utilize the SwiGLU activation~\citep{swiglu}, while rotary positional embeddings (RoPE)~\citep{RoPe} are applied in self-attention modules with a parameter of $\theta=500000$~\citep{xiong2024effective}. For normalization, we implement RMSNorm \citep{RMSNorm}. Throughout all self-attention blocks employing conventional attention masking—such as \textit{block causal} and \textit{fixed-window block causal}—we leverage Flash attention~\citep{dao2022flashattention}, fixing the attention window at 512 tokens for masks of constant width. When dealing with cross-attention components—where masks are computed dynamically and conditioned on specific input patches—we instead rely on Flex Attention\footnote{\url{https://pytorch.org/blog/flexattention}}. This choice enables efficient, fused kernel implementations that markedly enhance training speed.

\subsection{{\textsc BLT}-Specific Hyperparameters}

In order to evaluate the performance of {\textsc BLT} models, our experimental analysis proceeds along two main lines: assessing scaling behavior and examining downstream task performance, considering models with sizes of 400M, 1B, 2B, 4B, and 8B. The associated architectural hyperparameters for these configurations are provided in Appendix~Table~\ref{tab:arch}.

For all {\textsc BLT} variants, the queries for the first cross-attention module in the local encoder are initialized via max-pooling. Hashing parameters are kept consistent, with $500,000$ hashes generated by a single hash function and n-gram lengths set between 3 and 8. The learning rate is fixed at $4e-4$ throughout. This particular learning rate was selected based on a hyperparameter sweep—spanning from $1e-3$ to $1e-4$—at the 400M and 1B parameter model sizes, where both token and {\textsc BLT} models showed optimal results with the same value.

Batch sizes used for scaling trend experiments on the Llama-2 corpus follow the recommendations outlined in~\cite{dubey2024llama}, or are scaled equivalently by byte size. Training optimization leverages the AdamW optimizer~\citep{adamw}, configured with $\beta_1=0.9$, $\beta_2=0.95$, and $\epsilon=10^{-8}$. We employ a cosine learning rate decay to zero, preceded by a linear warm-up phase over 2000 steps. Additionally, weight decay is set to 0.1, and global gradient norms are clipped at 1.0.

\section{Scaling Trends}
\label{section:scaling}

We provide a comprehensive overview of how byte-level models scale, with the intention of guiding the continued development of {\textsc BLT} models. Our analysis seeks to overcome the shortcomings observed in prior work on byte-level architectures by taking several new steps: (a) We analyze scaling dynamics under compute-optimal training conditions; (b) We build and assess 8B parameter models using substantial datasets (as much as 1T tokens or 4T bytes), and examine their performance across a range of downstream applications; and (c) We explore scaling properties in contexts where inference cost is controlled. Later in this paper, we further explore particular strengths associated with the modeling of byte-sequences.

\subsection{Parameter Matched Compute Optimal Scaling Trends}
\label{section:parameter-matched-scaling}
We employ the Llama 2 dataset to train several \textit{compute-optimal} bpe and {\textsc BLT} models, spanning four parameter scales from 1B to 8B. Subsequently, we visualize training flops alongside language modeling performance, using a representative portion of the training data mixture. The bpe models are optimized according to the model-to-data ratio recommended by Llama~3~\citep{dubey2024llama}, which designates the \textit{compute-optimal} regime for maximizing training data performance within a prescribed computational budget~\citep{hoffmann2022training}. This approach establishes a strong baseline for subsequent comparison. To facilitate fair evaluation, each bpe model has a corresponding {\textsc BLT} model trained on identical data, with a Latent Transformer configured to have the same size and architecture as the corresponding bpe Transformer.

As shown in~\autoref{fig:scaling-compute-opt} (right), {\textsc BLT} models consistently equal or surpass the performance of bpe baselines, a pattern that persists with increasing model size and computational resources. To our knowledge, {\textsc BLT} represents the first byte-level Transformer design to demonstrate comparable scaling behavior to BPE-based systems when operating in compute-optimal settings. This supports the hypothesis that the optimal parameter-to-compute ratio for bpe models is similarly applicable to {\textsc BLT}, or at minimum, does not differ significantly.

Both advances in model architecture and the introduction of dynamic patching mechanisms are essential for achieving bpe scaling parity. In ~\autoref{fig:scaling-compute-opt} (left), we present a comparison between models utilizing space-patching and the Llama 3 architecture. For this purpose, SpaceByte \citep{slagle2024spacebyte} is approximated by implementing {\textsc BLT} space-patching, omitting n-gram embeddings and cross-attention. While SpaceByte surpasses Megabyte in terms of performance, there remains a significant gap relative to Llama 3. \autoref{fig:scaling-compute-opt} (right) demonstrates that gains are realized when both architectural enhancements and dynamic patching are combined. At scale, {\textsc BLT} models achieve competitive results with leading tokenizer-based architectures such as Llama 3.

Additionally, we note that model performance for tokenizer-based approaches is influenced by tokenizer selection; models utilizing the Llama-3 tokenizer show superior results compared to those employing the Llama-2 tokenizer on equivalent datasets.

In summary, our {\textsc BLT} architecture occupies an intermediate position between Llama 2 and 3 when evaluated with substantially larger patch sizes. While the Llama 2 and 3 bpe tokenizers produce tokens averaging 3.7 and 4.4 bytes respectively, {\textsc BLT} can follow analogous scaling behavior when utilizing patch sizes of 6 or even 8 bytes on average. Since inference flops decrease as the average patch size increases, employing 8-byte patches can result in almost 50\% reduction in inference computational cost. Moreover, models utilizing larger patches tend to improve in performance as both the model scale and dataset size increase. Notably, although the 8-byte patch size configuration of {\textsc BLT} initially underperforms relative to bpe Llama 2 at 1B scale, it surpasses bpe performance at the 7B scale. This pattern indicates that models could benefit even further from larger patch sizes as scaling continues and resources expand, potentially making even larger patch sizes viable in the future.

\subsection{Beyond Compute Optimal Task Evaluations}
\label{section:evals}
To further investigate scaling behavior, we extend training of an 8B {\textsc BLT} model past the compute-optimal threshold using the {\textsc BLT}-1T dataset, which is both larger and of superior quality, and we assess its performance using a diverse array of established classification and generation benchmarks. For our evaluation, we consider tasks in common sense reasoning, world knowledge, and code generation:
\paragraph{Classification tasks} encompass ARC-Easy (0-shot)~\citep{Eval_ARC}, Arc-Challenge (0-shot)~\citep{Eval_ARC}, HellaSwag (0-shot)~\citep{Eval_hellaswag}, PIQA (0-shot)~\citep{Eval_piqa}, and MMLU (5-shot)~\citep{hendrycks2020measuring}. We utilize a prompt-based scoring technique, where we compute the likelihood over the answer choices, and present the mean accuracy.
\paragraph{Coding related generation tasks:} We evaluate performance using pass@1 metrics on MBPP (3-shot)~\citep{Eval_MBPP} and HumanEval (0-shot)~\citep{Eval_HumanEval}, thereby testing the model’s capacity for Python code synthesis.

In \autoref{tab:evals}, we present a comparative analysis of three different models that are all trained on the {\textsc BLT}-1T dataset: a bpe-based model utilizing the Llama 3 tokenizer,\footnote{The Llama 3 tokenizer, featuring a 128k vocabulary, is selected due to its superior performance over the Llama 2 tokenizer with a 32k vocabulary.} and two {\textsc BLT} model variants. The first {\textsc BLT} variant incorporates a space-patching technique ({\textsc BLT}-Space), while the second applies an entropy-driven patching mechanism ({\textsc BLT}-Entropy). The latter includes an approximate monotonicity constraint and reinitializes the context for the entropy-based model at every new line, following the methodology outlined in~\autoref{section:entropy-context}. Training for all three approaches is carried out under an equivalent flop budget. Notably, for {\textsc BLT}-Entropy, the entropy threshold at inference is reduced from 0.6 to 0.1, which is observed to yield improved task accuracy, albeit with an associated increase in the number of inference steps required.

The {\textsc BLT}-Entropy model demonstrates superior performance compared to the Llama 3 model in 4 out of 7 evaluated tasks, despite both models being trained on an equal volume of data. This advantage is likely attributable to two primary factors: (1) enhanced utilization of computational resources through dynamic patching, and (2) the explicit modeling of information at the byte level rather than at the token level.

Conversely, {\textsc BLT}-Space exhibits inferior performance compared to the Llama 3 tokenizer on all tasks except one; nevertheless, it yields a notable decrease in inference flops due to its greater average patch size of 6 bytes. In contrast, both the bpe and entropy-patching approaches generate an average patch size near 4.5 bytes within the training data mix. Given an identical training budget, the model employing the larger patch size processes 30\% more data than the other two approaches, which could cause {\textsc BLT} to diverge further from the compute-optimal regime.

\subsection{Patches Scale Better Than Tokens} Utilizing {\textsc BLT} models allows for the concurrent scaling of both the model and the patch size without exceeding the available flop budget for training and inference, and with the training dataset size remaining unchanged. Patch-based architectures uniquely permit indefinite enlargement of the patch size, surpassing the limitations inherent to fixed-vocabulary token-based models as outlined in Section~\ref{section:bpe-patch}. Increasing patch lengths reduces computational requirements, enabling those resources to be redirected towards expanding the global latent transformer, since it needs to operate less frequently.

To evaluate whether increasing model size while decreasing the number of steps on larger patches yields improved performance compared to smaller models executing more steps, we perform a fixed inference scaling analysis. We begin with 400m and 3.6B parameter models utilizing the Llama 2 tokenizer, and then identify flop-equivalent counterparts using the Llama 3 tokenizer, along with {\textsc BLT}-Entropy models featuring mean patch sizes of 6 and 8 bytes on the training datamix (refer to~\autoref{tab:fixed-inf-params} for further specifications). For the models with patch size 8, we employ 3 encoder layers, as opposed to a single layer for others. Each configuration is trained across a range of training flop allocations.

As depicted in \autoref{fig:fixed_inference_scaling}, {\textsc BLT} models demonstrate more favorable scaling behavior compared to tokenization-based models across both categories of inference flop requirements. Although BPE models outperform in scenarios with limited training resources, their advantage diminishes as {\textsc BLT} models rapidly overtake them shortly after surpassing the compute-optimal point. This suggests that, in many cases, investing additional resources during the initial pretraining phase can be advantageous, resulting in superior model performance when inference computation is held constant. A clear illustration of this is seen with the 8B-class models, such as Llama 3.1, which have been pretrained on data volumes two orders of magnitude greater than those optimal from a compute perspective for models of that scale.

The threshold at which {\textsc BLT} starts to outperform token-based models has moved marginally closer to the compute-optimal configuration with the adoption of models from a higher flop class, shifting from approximately 3x to 2.5x the compute-optimal allocation. In parallel, the model using patch size 8 in this larger flop regime demonstrates a more accelerated improvement, surpassing the performance of alternative models at an earlier stage. As outlined in Section~\ref{section:parameter-matched-scaling}, when scaling up model size, increased patch sizes tend to achieve results more comparable to BPE models. This behavior can partly be explained by the reduced proportion of overall flops devoted to the byte-level Encoder and Decoder components, which do not scale as rapidly as the Latent Transformer. Specifically, when the total parameter count is increased twentyfold, from 400M to 8B, the local parameters within {\textsc BLT} only expand by a factor of around two. This is a crucial observation because increases in patch size exclusively influence the computational demand of the patch Latent Transformer, leaving the byte-level components largely unaffected. Consequently, the proportion of {\textsc BLT}-Entropy with ps=8 compared to Llama 2's model size increased slightly, from 1.6x to 1.7x, as the scale of the model was enlarged.

In conclusion, our investigation into patch-length scaling reveals that the {\textsc BLT} patch-based model architecture exhibits superior scaling behavior when both the patch size and the overall model size are increased together. These scaling advantages appear to not only persist but may also be further enhanced as model size grows.

\section{Byte Modeling Improves Robustness} We further assess the robustness of {\textsc BLT} in comparison to token-based models that are deprived of direct byte-level representation, and we introduce a method for augmenting pretrained token-based models with byte-level capabilities.
\label{sec:robustness}

\subsection{Character-Level Tasks}

One of the initial driving reasons for developing byte-level models was their ability to handle noise present at the byte level in input data, along with their capacity to recognize the subcomponents that make up tokens—an area where traditional tokenizer-based architectures typically encounter difficulties. To systematically assess these characteristics, we conducted supplementary experiments using benchmark datasets designed to test both the resilience of models against noisy input and their sensitivity to character-level structure, across English and various other languages, as well as for digits and phonemes. These experimental findings are summarized in Table~\ref{tab:char_tasks}.

\paragraph{Noisy Data} To evaluate and contrast the resilience of tokenizer-based models and {\textsc BLT}, we generate noisy variants of the benchmark classification datasets referenced in Section~\ref{section:evals}. We utilize five separate character-level noising methodologies to introduce textual perturbations: (a) \textit{AntSpeak}: Rewrites the text such that each character is capitalized and separated by spaces. (b) \textit{Drop}: Eliminates 10\% of the text's characters at random. (c) \textit{RandomCase}: Randomly assigns uppercase or lowercase status to 50\% of characters each. (d) \textit{Repeat}: Selects 20\% of the characters at random and duplicates them up to four times each. (e) \textit{UpperCase}: Converts the entirety of the text to uppercase letters. For the purpose of evaluation, each of these noising operations is applied independently to either the prompt, the completion, or both; scores are averaged over these configurations. Results are shown in Table \ref{tab:char_tasks} for noised HellaSwag~\citep{Eval_hellaswag}, revealing that {\textsc BLT} consistently surpasses tokenizer-based baselines in robustness. On average, it outperforms a model trained on identical data by 8 points and further achieves better results than the Llama 3.1 model, which is trained on a significantly larger corpus.

\paragraph{Phonology - Grapheme-to-Phoneme (G2P)} We evaluate {\textsc BLT}'s performance in converting sequences of graphemes (letters that constitute a word) into their corresponding phonemic transcriptions (pronunciations). As shown in Table \ref{tab:char_tasks}, we report outcomes for the G2P task under a 5-shot learning regime using Phonology Bench~\citep{suvarna2024phonologybench}. The results demonstrate that {\textsc BLT} surpasses the Llama 3 1T tokenizer-based model baseline on this evaluation.

\paragraph{CUTE}
To evaluate character-level comprehension, we test {\textsc BLT} on the CUTE benchmark~\citep{edman2024cute}, which includes a suite of tasks generally divided into three groups: compositional understanding, orthographic similarity assessment, and sequence manipulation abilities. These tasks present a notable challenge for most models reliant on tokenization, which tend to have internalized token spellings but often fail to leverage this knowledge effectively for manipulating textual content. As indicated in Table~\ref{tab:char_tasks}, {\textsc BLT}-Entropy achieves a margin exceeding 25 points over both BPE Llama 3 variants on this evaluation. Particularly, the model displays outstanding accuracy in character manipulation, attaining a score of 99.9\% on both spelling-related tasks. Such substantial gains—achieved even though {\textsc BLT} was trained on sixteen times less data compared to Llama 3.1—highlight the inherent difficulty BPE models encounter in acquiring character-level proficiency. Figure \ref{fig:cute_examples} showcases representative instances where Llama 3's tokenizer-based approach falters but our {\textsc BLT} model excels. Notably, word deletion and insertion remain the only categories where BPE outperforms. While such manipulations at the word level pose challenges for a byte-level architecture, the difference in performance is not pronounced, and it may be simpler to construct word-level understanding from characters than to infer character-level details from word-based representations. Across all tasks, the evaluation methodology and original Huggingface prompts remain consistent. Further, additional prompt engineering could potentially enhance performance for BPE-based models.

\paragraph{Low Resource Machine Translation} We assess {\textsc BLT} on both directions of translation across twenty-one low-resource languages spanning six major language families and incorporating a range of scripts, leveraging the FLORES-101 benchmark~\citep{goyal2022flores}. SentencePiece BLEU scores are presented in Table~\ref{tab:flores}. The findings indicate that {\textsc BLT} surpasses a counterpart utilizing the Llama 3 tokenizer, delivering a 2-point overall gain for translation into English and a 0.5-point improvement for translation from English. While performance on widely studied language pairs is on par with or marginally superior to Llama 3, {\textsc BLT} exhibits notable improvements over Llama 3 in many pairs involving languages from low-resource families. This highlights the capacity of byte-level modeling to effectively generalize across diverse, less common byte sequences.

\subsection{Training {\textsc BLT} from Llama 3}
\label{sec:distillation}
We investigate a strategy in which {\textsc BLT} models take advantage of pre-existing, pre-trained models that utilize tokenizers to enhance and expedite the training process. Specifically, we initialize the global transformer parameters of {\textsc BLT} using those derived from a pre-trained Llama 3.1 model. Following this initialization, we continue updating the global transformer’s weights with a learning rate set to one-tenth of that used for both the local encoder and decoder, applying the same number of steps as used in Llama 3’s training procedure. Table~\ref{tab:distill} offers a direct comparison against a standard {\textsc BLT} baseline. The results indicate that initializing {\textsc BLT} from Llama 3.1 leads to markedly better outcomes than both the original Llama 3 and the default {\textsc BLT} models trained for equivalent FLOPs. Furthermore, when benchmarked against the {\textsc BLT}-Entropy variant in Table~\ref{tab:evals}—which benefitted from a much more extensive training dataset (1T tokens)—the {\textsc BLT} initialized from Llama 3.1 still registers higher MMLU scores. This suggests that such a distillation-based approach can considerably diminish the computational resources required for training without compromising performance.

This approach can alternatively be interpreted as adapting tokenizer-based architectures to function without tokenizers, thereby enabling the transformation of a pre-trained LLaMA 3.1 model into a {\textsc BLT} framework. For a thorough evaluation, we present results for the original LLaMA 3.1 model—trained on 15T tokens—in Table \ref{tab:distill}, and compare its outcomes to those of the LLaMA 3-based {\textsc BLT}. While the transition leads to only modest reductions in performance on MMLU and HumanEval benchmarks, more pronounced decreases are observed across the remaining tasks. These findings indicate that additional advancements are necessary to maximize the benefits of the pre-trained architecture, such as refining the selection of training data and fine-tuning hyperparameter configurations.

\section{Ablations and Discussion}
\label{section:ablations}
Here, we present ablation studies to motivate design decisions for the {\textsc BLT} architecture, including the specifics of the patching strategy and selected hyper-parameters for the 8B parameter {\textsc BLT} model trained on the {\textsc BLT}-1T dataset.

\paragraph{Entropy Model Hyper-parameters} To evaluate how adjustments in the entropy model's size and the length of the context window influence scaling behavior, we train a series of byte-level entropy transformer models, varying the parameter count from 1 million to 100 million and the context window from 64 up to 512 tokens. The relationship between bits-per-byte (bpb) and training flop across scales is illustrated for our $400m$ and $1b$ {\textsc BLT} models trained on the Llama-2 dataset in \autoref{fig:entropy_ablation}. Our observations indicate that both increasing the entropy model's size and context window length enhances scaling efficiency, although these benefits begin to taper off past the 50 million parameter mark.

\paragraph{Types of Patching}

We conduct an ablation study of the four patching methods outlined in Section~\ref{section:patching}: (1) Strided Patching utilizing strides of 4 and 6, (2) Patching determined by whitespace segmentation, (3) Patch creation guided by the Llama 3 tokenizer through BPE Tokenizer patching, and (4) An entropy-driven patching approach leveraging a lightweight byte-level language model.

Although dynamic patching shortens the actual sequence lengths, we ensure that all patching approaches have matched sequence lengths so that the available context remains consistent. Across both training and inference, every model processes an equivalent expected number of bytes per sequence, which controls for the ability to utilize larger context windows and eliminates related confounders. The results of these ablation studies are presented in Figure~\ref{fig:scaling-compute-opt}. Notably, all alternative patching strategies yield better performance than static patching; among these, space patching demonstrates performance nearly on par with dynamic entropy based patching.

\autoref{tab:evals_ablation} reports benchmark results for {\textsc BLT} models, evaluating the performance of tokenizer-based approaches, space patching, and entropy-based patching, each trained for an optimal duration on the Llama 2 dataset~\citep{dubey2024llama}. While space patching represents a more straightforward method that avoids the need to deploy an entropy model dynamically during training, our findings indicate that the improvements gained with entropy-based patching—previously identified in scaling experiments (Section~\ref{section:scaling})—are retained on a variety of downstream benchmark tasks.\footnote{Space patching results reflect earlier experiments lacking cross-attention, yet comparable trends persist when cross-attention is included.}

\paragraph{Cross-Attention} 
In~\autoref{tab:crossattn_ablations_1b}, we present an ablation study assessing the incorporation of cross-attention at different locations within the encoder and decoder of {\textsc BLT}. For encoder cross-attention, we investigate three initialization strategies for the queries: (1) a single, shared learned embedding across all global states, (2) a hash-based embedding derived from the bytes comprising each patch, and (3) a pooled version of the encoder's hidden representation of the patch bytes at the corresponding encoder layer.

Our results indicate that implementing cross-attention within the \textit{decoder} yields the highest effectiveness. When applied to the encoder, cross-attention offers a modest benefit, and this is observed solely when query initialization is performed through pooling. Furthermore, cross-attention proves particularly beneficial for Common-Crawl data, with the most pronounced gains occurring for larger patch sizes.

\paragraph{n-gram Hash Embeddings} 

We experiment with n-gram hash embedding vocabulary sizes of 0, 100K, 200K, and 400K, and display the outcomes in Table~\ref{tab:ngram_ablations_1b}. The analysis shows that hash embeddings universally benefit performance across domains, with pronounced advantages on Wikipedia and Github (yielding a 0.04 bpb improvement versus a 0.01 bpb improvement at 8B parameters after 15k training steps). When scaling to 8B parameters, reducing the number of hashes from 500K to 300K produces a marginal performance change of approximately 0.001 bpb after 15k steps. These findings suggest that hash embeddings are crucial for enabling {\textsc BLT} to achieve parity with models that use tokenizers; nonetheless, returns significantly taper off beyond 300K hashes. Furthermore, our results suggest that performance improvements gained from hashes are mostly additive to those from cross-attention, as each provides enhancements on separate sets of data.

\paragraph{Local Model Hyperparameters} Table \ref{tab:local_layers} presents an ablation of different configurations for the number of layers utilized in both the local encoder and decoder. The results indicate that {\textsc BLT} achieves strong performance with hash n-gram embeddings when the encoder is kept minimal—consisting of only a single layer—while employing a comparatively more complex, multi-layer decoder.

\section{Related Work}
\label{section:related_work}

\textbf{Character-Level RNNs:} The task of Character Language Modeling has long been central to neural network research~\citep{sutskever2011generating,mikolov2012subword,graves2013generating}, due to its inherent capacity to naturally handle out-of-vocabulary terms, eliminating the reliance on traditional back-off strategies. \cite{kim2016character} developed a model exclusively processing character-level input, combining convolutional and highway networks preceding LSTM-based RNNs, and succeeded in reaching parity with then state-of-the-art RNN language models on English, while surpassing their performance on languages with complex morphology—a frequently highlighted benefit of character-based LLMs. Similarly, \cite{kenter2018byte} employed byte-level LSTM architectures for machine comprehension, which again outperformed word-level approaches on morphologically intricate languages like Turkish and Russian. In related research, \cite{zhang2015character} implemented convolutional models based on characters for text classification, achieving superior results over word-level baselines on specific tasks. The hierarchical LSTM framework proposed by \citet{chung2019hierarchical}, integrating boundary detectors at each layer, aimed to automatically uncover latent textual hierarchies and thereby enhance the effectiveness of character-level language modeling. Additionally, ByteNet by \cite{kalchbrenner2016neural} applied convolutional layers to character input—contrasting with attention mechanisms—for machine translation tasks.

\textbf{Character-Level Transformers:} The introduction of attention-based transformer architectures~\citep{vaswani2017attention}, alongside techniques such as subword tokenization~\citep{sennrich-etal-2016-neural}, led to notable advancements in neural language modeling and benchmark performances. Nonetheless, the choice of word and sub-word representations inherently establishes a certain inductive bias dictating the abstraction level that models process. Seeking to merge the effectiveness of transformers with encouraging trends in character-level language modeling, \citet{al2019character} explored training extremely deep transformer models and leveraged auxiliary loss functions to construct transformer-based approaches that surpassed previous LSTM-driven character language models. Despite these improvements, their models continued to underperform in comparison to word-level LLMs. Moreover, GPT-2~\citep{radfordgpt2} reported that when working with extensive datasets, such as the 1 Billion Word Benchmark, language models operating at the byte level still lagged behind their word-level counterparts in terms of competitiveness.

\cite{Choe2019BridgingTG} showed that transformer-based language models operating at the byte level can surpass subword-level LLMs with equivalent parameter counts; however, these byte-level architectures are significantly more resource-intensive and require considerably longer training times. In a related vein, \cite{el2020characterbert} introduced CharFormer, a BERT-based model that constructs word representations via convolutional operations on character embeddings and demonstrates domain-specific gains in the medical field, but again at the cost of increased computational demands. The CANINE model, presented by \cite{clark2022canine}, consists of an encoder-only architecture with 150M parameters that processes input at the character level. This architecture incorporates a deep transformer core reminiscent of the structure in our global model, supplemented by a local transformer and strided convolutions to perform downsampling on the character sequences. CANINE surpasses the token-level encoder counterpart (mBERT) when evaluated on multilingual downstream tasks. ByT5~\citep{xue2022byt5} investigated strategies for building encoder-decoder models that process inputs at the byte level, specifically avoiding patching mechanisms. Although ByT5 displays notable resilience to input noise and achieves similar performance to token-based models even with only one-quarter as much data, the absence of patching necessitates computationally costly attention over each byte, making the process especially demanding. The use of bytes as modeling units, rather than subwords, greatly lengthens input sequences, complicating efficient scaling of byte-level models. Recently, leveraging the Mamba Architecture~\citep{gu2023mamba}—which provides the capability of fixed-size memory management over long contexts—\cite{wang2024mambabyte} trained a byte-level Mamba model without patching and demonstrated that it outperforms byte-level transformer architectures in flop-matched experiments at the 350M parameter scale, specifically in terms of bits-per-byte across various datasets.

\textbf{Patching-based approaches:} Leveraging patching techniques can significantly reduce the excessive computational cost, measured in flops, often associated with byte-level LLMs, without necessarily compromising model quality. Several studies have reported early successes at limited model sizes and training data volumes. For instance, \cite{nawrot-etal-2022-hierarchical} explored approaches employing fixed patching for both downsampling and upsampling, introducing the hourglass transformer, which surpassed alternative byte-level methods at the 150M parameter scale. Building on this, \cite{nawrot-etal-2023-efficient} advanced dynamic patching strategies. Their work introduced various boundary-predictor mechanisms—one learned in a fully end-to-end framework, another trained with explicit tokenizer supervision, and a third utilizing entropy-based segmentation akin to {\textsc BLT}. Their empirical results indicate that such methodologies surpassed contemporaneous vanilla transformers in language modeling benchmarks when trained with approximately 40M parameters on a corpus of about 400M tokens. Additionally, \cite{lester2024training} studied training on sequences compressed via arithmetic coding, which allowed for compression levels exceeding those achievable with BPE. Through the application of an equal-info windows method, their models exceeded the performance of byte-level baselines on language modeling tasks, although they still trailed behind models using subword tokenization strategies.

Our research is heavily influenced by and shares the closest affinity with MegaByte~\citep{yu2023megabyte}, a causal decoder-only LLM that leverages a fixed, static patching mechanism paired with representation concatenation to transition from bytes to patches. A local model within the decoder is then used to translate these patches back to bytes. Their experiments indicate that, at a 1B parameter scale and on datasets encompassing approximately 400B bytes, MegaByte achieves performance comparable to models utilizing tokenizers. Throughout our experiments, we systematically evaluate MegaByte and observe that static patching generally underperforms compared to state-of-the-art, compute-efficient tokenizer-based models under flop-normalized conditions. Our findings further reveal that {\textsc BLT} substantially narrows this gap. In parallel, \cite{slagle2024spacebyte} report a similar lag with MegaByte, and advocate for enhancing the static patching approach by incorporating patches aligned with whitespaces and similar bytes, as well as integrating a local encoder model. Their results display performance gains over token-based transformer models within a compute-normalized context across domains like Github and arXiv at the 1B parameter level. We extend this line of investigation with experiments on their model and demonstrate that to truly rival state-of-the-art token-based systems such as Llama 3, further advances in byte-level model architecture are essential as scale increases.

\section{Limitations and Future Work}

In this study, we utilize the optimal training steps established for Llama~3~\citep{dubey2024llama} to guide our architectural decisions. It is important to note, though, that these scaling laws were originally derived from BPE-level transformer models and might not yield the most effective data-to-parameter size ratios for the {\textsc BLT} framework. Deriving scaling laws specifically tailored to {\textsc BLT}, which could result in superior scaling characteristics for this architecture, remains an area for future investigation. Moreover, the majority of experiments herein have been performed at scales up to 1B parameters; it remains an open question whether the most advantageous architectural configurations will persist as model sizes extend to 8B parameters and above, where further performance enhancements could become apparent.

Current transformer libraries and repositories are primarily optimized for the efficiency of tokenizer-based transformer models. Although we provide theoretical flop-equivalent experiments and leverage optimized solutions like FlexAttention for components that differ from standard transformer designs, our implementations might still lag behind tokenizer-based counterparts in terms of wall-clock performance and could potentially see improvements from additional optimization efforts.

Although {\textsc BLT} currently employs an entropy model for patching that is trained independently, an appealing avenue for future research is to explore end-to-end training of the patching model itself. Section \ref{sec:distillation} outlines preliminary studies that suggest it is feasible to "byte-ify" models with token-based tokenizers, such as Llama 3, which have been trained on corpora exceeding 10T tokens. This is achieved by initializing and freezing the parameters of the global transformer using the pretrained weights. Continued investigation in this area may reveal techniques that not only maintain the advantages conferred by bytefying but also enable superior performance compared to existing tokenizer-based architectures—potentially without the necessity of training these systems anew.

\section{Conclusion}
\label{section:conclusion}

In this work, we introduce the Byte Latent Transformer (\textbf{BLT}), a novel framework that fundamentally revisits the standard reliance on predetermined token vocabularies in large language models. By adopting a flexible and trainable scheme for assembling bytes into variable-length patches, {\textsc BLT} dynamically directs computational focus according to the intricacy of the input, thereby achieving notable gains in both computational efficiency and resilience. Through a comprehensive scaling analysis, we find that {\textsc BLT} models attain parity with established tokenization-dependent models such as Llama 3 at scales extending up to 8B parameters and 4T bytes. Moreover, they permit up to 50\% reduction in inference flops for marginal sacrifices in benchmark performance. Importantly, {\textsc BLT} introduces a previously unavailable scaling avenue: model and patch size can be expanded jointly while maintaining a fixed computation budget during inference—an asset particularly beneficial under real-world compute constraints. Additionally, by processing unsegmented byte inputs, {\textsc BLT} enhances the capability of the model to process rare and outlier data, achieving greater resistance to noisy examples and yielding more nuanced modeling of sub-word patterns. Collectively, these findings identify {\textsc BLT} as a compelling alternative to legacy tokenization strategies, offering an efficient, adaptable, and robust methodology for next-generation language modeling.

\end{document}